% GENERATED FILE. Edit canonical lessons, then run npm run generate:books.
% canonical-source-hash: fnv1a64:b1d10dd1c44837e2
% canonical-lessons: GU-C545-asan, GU-C545-jatil, GU-C545-mahattva, GU-C545-ajmayas, GU-C545-yukti

\chapter{Easy or Complicated}
\label{ch:gu-easy-or-complicated}
\hlchaptermodality{545}

\begin{chapteropening}
\textbf{By the end of this chapter:} \emph{I can say whether a way of doing something is easy or complicated.}

Complete the last lesson of chapter 545: I can say whether a way of doing something is easy or complicated.
\end{chapteropening}

\section[\texorpdfstring{\gu{આસાન} (āsān)}{આસાન (āsān)}]{\gu{આસાન} (āsān) — effortless}

\label{lesson:GU-C545-asan}



\begin{quote}
Before the new one: say the Gujarati for a copy, then the Gujarati for evenly matched.
\end{quote}

\subsection*{You'll want to know: \gu{આસાન}}
\textbf{\gu{આસાન}} — \emph{āsān} — \textquotedblleft{}effortless\textquotedblright{}.

\begin{quote}
\gu{જો} \gu{તમે} \gu{આ} \gu{રીત} \gu{વાપરશો} \gu{તો} \gu{કામ} \gu{આસાન} \gu{બનશે}.
\end{quote}

\begin{quote}
\emph{jo tame ā rīt vāparśo to kām āsān banśe.}
\end{quote}

\begin{quote}
If you use this method, the work will become easy.
\end{quote}

\begin{cousinweb}
\textbf{\gu{આસાન}} (\emph{āsān}) is usually traced to Persian \emph{āsān}, \textquotedblleft{}easy\textquotedblright{}; the everyday Gujarati word beside it is \textbf{\gu{સહેલું}}.
\end{cousinweb}

\begin{grammarlens}[title={What you've built}]
One more word for saying what happens if.
\end{grammarlens}

\subsection*{Your turn}
\begin{itemize}
\raggedright
  \item \emph{Say it:} \emph{āsān}
  \item \emph{Say it:} \emph{āsān}, once more
  \item \emph{Say it:} \emph{jo tame ā rīt vāparśo to kām āsān banśe.} — If you use this method, the work will become easy.
  \item \emph{Recall:} say the Gujarati for a copy, then the Gujarati for evenly matched, then say \emph{āsān} again
\end{itemize}

\begin{tcolorbox}[breakable,colback=teal!4,colframe=teal!35!black,title={Before you move on}]
What does \gu{આસાન} mean? (\textquotedblleft{}Effortless\textquotedblright{}.) Say it once more.
\end{tcolorbox}
\section[\texorpdfstring{\gu{જટિલ} (jaṭil)}{જટિલ (jaṭil)}]{\gu{જટિલ} (jaṭil) — complicated}

\label{lesson:GU-C545-jatil}



\begin{quote}
Before the new one: say the Gujarati for evenly matched, then the Gujarati for effortless.
\end{quote}

\subsection*{You'll want to know: \gu{જટિલ}}
\textbf{\gu{જટિલ}} — \emph{jaṭil} — \textquotedblleft{}complicated\textquotedblright{}.

\begin{quote}
\gu{જો} \gu{નિયમો} \gu{જટિલ} \gu{હશે} \gu{તો} \gu{કોઈ} \gu{તે} \gu{નહીં} \gu{સમજે}.
\end{quote}

\begin{quote}
\emph{jo niymo jaṭil haśe to koī te nahī̃ samje.}
\end{quote}

\begin{quote}
If the rules are complicated, nobody will understand them.
\end{quote}

\begin{cousinweb}
\textbf{\gu{જટિલ}} (\emph{jaṭil}) is a tatsama, taken from Sanskrit \emph{jaṭila}, \textquotedblleft{}wearing matted hair\textquotedblright{} (Monier-Williams): tangled like matted locks.
\end{cousinweb}

\begin{grammarlens}[title={What you've built}]
One more word for saying what happens if.
\end{grammarlens}

\subsection*{Your turn}
\begin{itemize}
\raggedright
  \item \emph{Say it:} \emph{jaṭil}
  \item \emph{Say it:} \emph{jaṭil}, once more
  \item \emph{Say it:} \emph{jo niymo jaṭil haśe to koī te nahī̃ samje.} — If the rules are complicated, nobody will understand them.
  \item \emph{Recall:} say the Gujarati for evenly matched, then the Gujarati for effortless, then say \emph{jaṭil} again
\end{itemize}

\begin{tcolorbox}[breakable,colback=teal!4,colframe=teal!35!black,title={Before you move on}]
What does \gu{જટિલ} mean? (\textquotedblleft{}Complicated\textquotedblright{}.) Say it once more.
\end{tcolorbox}
\section[\texorpdfstring{\gu{મહત્ત્વ} (mahattva)}{મહત્ત્વ (mahattva)}]{\gu{મહત્ત્વ} (mahattva) — importance}

\label{lesson:GU-C545-mahattva}



\begin{quote}
Before the new one: say the Gujarati for effortless, then the Gujarati for complicated.
\end{quote}

\subsection*{You'll want to know: \gu{મહત્ત્વ}}
\textbf{\gu{મહત્ત્વ}} — \emph{mahattva} — \textquotedblleft{}importance\textquotedblright{}.

\begin{quote}
\gu{જો} \gu{તમે} \gu{સમયનું} \gu{મહત્ત્વ} \gu{સમજશો} \gu{તો} \gu{તમે} \gu{ક્યારેય} \gu{મોડા} \gu{નહીં} \gu{પડો}.
\end{quote}

\begin{quote}
\emph{jo tame samaynũ mahattva samajśo to tame kyārey moḍā nahī̃ paḍo.}
\end{quote}

\begin{quote}
If you understand the value of time, you will never be late.
\end{quote}

A neuter noun.

\begin{cousinweb}
\textbf{\gu{મહત્ત્વ}} (\emph{mahattva}) is a tatsama, taken from Sanskrit \emph{mahattva}, \textquotedblleft{}greatness, magnitude\textquotedblright{} (Monier-Williams).
\end{cousinweb}

\begin{grammarlens}[title={What you've built}]
One more word for saying what happens if.
\end{grammarlens}

\subsection*{Your turn}
\begin{itemize}
\raggedright
  \item \emph{Say it:} \emph{mahattva}
  \item \emph{Say it:} \emph{mahattva}, once more
  \item \emph{Say it:} \emph{jo tame samaynũ mahattva samajśo to tame kyārey moḍā nahī̃ paḍo.} — If you understand the value of time, you will never be late.
  \item \emph{Recall:} say the Gujarati for effortless, then the Gujarati for complicated, then say \emph{mahattva} again
\end{itemize}

\begin{tcolorbox}[breakable,colback=teal!4,colframe=teal!35!black,title={Before you move on}]
What does \gu{મહત્ત્વ} mean? (\textquotedblleft{}Importance\textquotedblright{}.) Say it once more.
\end{tcolorbox}
\section[\texorpdfstring{\gu{અજમાયશ} (ajmāyaś)}{અજમાયશ (ajmāyaś)}]{\gu{અજમાયશ} (ajmāyaś) — a try-out}

\label{lesson:GU-C545-ajmayas}



\begin{quote}
Before the new one: say the Gujarati for complicated, then the Gujarati for importance.
\end{quote}

\subsection*{You'll want to know: \gu{અજમાયશ}}
\textbf{\gu{અજમાયશ}} — \emph{ajmāyaś} — \textquotedblleft{}a try-out\textquotedblright{}.

\begin{quote}
\gu{જો} \gu{અજમાયશમાં} \gu{મશીન} \gu{બરાબર} \gu{ચાલશે} \gu{તો} \gu{અમે} \gu{તે} \gu{ખરીદીશું}.
\end{quote}

\begin{quote}
\emph{jo ajmāyaśmā̃ maśīn barābar cālśe to ame te kharīdīśũ.}
\end{quote}

\begin{quote}
If the machine runs properly in the trial, we'll buy it.
\end{quote}

A feminine noun.

\begin{cousinweb}
\textbf{\gu{અજમાયશ}} (\emph{ajmāyaś}) is usually traced to Persian \emph{āzmāyish}, \textquotedblleft{}a trial, a test\textquotedblright{}; the verb \textbf{\gu{અજમાવવું}}, to try out, belongs with it.
\end{cousinweb}

\begin{grammarlens}[title={What you've built}]
One more word for saying what happens if.
\end{grammarlens}

\subsection*{Your turn}
\begin{itemize}
\raggedright
  \item \emph{Say it:} \emph{ajmāyaś}
  \item \emph{Say it:} \emph{ajmāyaś}, once more
  \item \emph{Say it:} \emph{jo ajmāyaśmā̃ maśīn barābar cālśe to ame te kharīdīśũ.} — If the machine runs properly in the trial, we'll buy it.
  \item \emph{Recall:} say the Gujarati for complicated, then the Gujarati for importance, then say \emph{ajmāyaś} again
\end{itemize}

\begin{tcolorbox}[breakable,colback=teal!4,colframe=teal!35!black,title={Before you move on}]
What does \gu{અજમાયશ} mean? (\textquotedblleft{}A try-out\textquotedblright{}.) Say it once more.
\end{tcolorbox}
\section[\texorpdfstring{\gu{યુક્તિ} (yukti)}{યુક્તિ (yukti)}]{\gu{યુક્તિ} (yukti) — a clever trick}

\label{lesson:GU-C545-yukti}



\begin{quote}
Before the new one: say the Gujarati for importance, then the Gujarati for a try-out.
\end{quote}

\subsection*{You'll want to know: \gu{યુક્તિ}}
\textbf{\gu{યુક્તિ}} — \emph{yukti} — \textquotedblleft{}a clever trick\textquotedblright{}.

\begin{quote}
\gu{જો} \gu{આ} \gu{યુક્તિ} \gu{કામ} \gu{કરશે} \gu{તો} \gu{આપણે} \gu{ઘણો} \gu{સમય} \gu{બચાવીશું}.
\end{quote}

\begin{quote}
\emph{jo ā yukti kām karśe to āpaṇe ghaṇo samay bacāvīśũ.}
\end{quote}

\begin{quote}
If this trick works, we'll save a lot of time.
\end{quote}

A feminine noun.

\begin{cousinweb}
\textbf{\gu{યુક્તિ}} (\emph{yukti}) is a tatsama, taken from Sanskrit \emph{yukti}, \textquotedblleft{}union, joining; a contrivance\textquotedblright{} (Monier-Williams).
\end{cousinweb}

\begin{grammarlens}[title={What you've built}]
One more word for saying what happens if.
\end{grammarlens}

\subsection*{Your turn}
\begin{itemize}
\raggedright
  \item \emph{Say it:} \emph{yukti}
  \item \emph{Say it:} \emph{yukti}, once more
  \item \emph{Say it:} \emph{jo ā yukti kām karśe to āpaṇe ghaṇo samay bacāvīśũ.} — If this trick works, we'll save a lot of time.
  \item \emph{Recall:} say the Gujarati for importance, then the Gujarati for a try-out, then say \emph{yukti} again
\end{itemize}

\begin{tcolorbox}[breakable,colback=teal!4,colframe=teal!35!black,title={Before you move on}]
What does \gu{યુક્તિ} mean? (\textquotedblleft{}A clever trick\textquotedblright{}.) Say it once more.
\end{tcolorbox}
